% GENERATED FILE. Edit canonical lessons, then run npm run generate:books.
% canonical-source-hash: fnv1a64:ef633c91a8c09b54
% canonical-lessons: TE-C293-jailu, TE-C293-masidu, TE-C293-carci, TE-C293-pancayati, TE-C293-ataka

\chapter{Jail, Mosque, Church, Village council, Loft}
\label{ch:te-jail-mosque-church-village-council-loft}
\hlchaptermodality{293}

\begin{chapteropening}
\textbf{By the end of this chapter:} \emph{I can say a jail, a mosque, a church, the village council and a loft.}

Complete the last lesson of chapter 293: I can say a jail, a mosque, a church, the village council and a loft.
\end{chapteropening}

\section[\texorpdfstring{\te{జైలు} (jailu)}{జైలు (jailu)}]{\te{జైలు} (jailu) — a jail}

\label{lesson:TE-C293-jailu}



\begin{quote}
Before the new one: say the Telugu for a fort, then the Telugu for a factory.
\end{quote}

\subsection*{You'll want to know: \te{జైలు}}
\textbf{\te{జైలు}} — \emph{jailu} — \textquotedblleft{}a jail\textquotedblright{}.

\begin{grammarlens}[title={What you've built}]
One more word for this chapter.
\end{grammarlens}

\subsection*{Your turn}
\begin{itemize}
\raggedright
  \item \emph{Say it:} \emph{jailu}
  \item \emph{Say it:} \emph{jailu}, once more
  \item \emph{Say it:} say \emph{jailu}
  \item \emph{Recall:} say the Telugu for a fort, then the Telugu for a factory, then say \emph{jailu} again
\end{itemize}

\begin{tcolorbox}[breakable,colback=teal!4,colframe=teal!35!black,title={Before you move on}]
What does \te{జైలు} mean? (\textquotedblleft{}A jail\textquotedblright{}.) Say it once more.
\end{tcolorbox}
\section[\texorpdfstring{\te{మసీదు} (masīdu)}{మసీదు (masīdu)}]{\te{మసీదు} (masīdu) — a mosque}

\label{lesson:TE-C293-masidu}



\begin{quote}
Before the new one: say the Telugu for a factory, then the Telugu for a jail.
\end{quote}

\subsection*{You'll want to know: \te{మసీదు}}
\textbf{\te{మసీదు}} — \emph{masīdu} — \textquotedblleft{}a mosque\textquotedblright{}.

\begin{grammarlens}[title={What you've built}]
One more word for this chapter.
\end{grammarlens}

\subsection*{Your turn}
\begin{itemize}
\raggedright
  \item \emph{Say it:} \emph{masīdu}
  \item \emph{Say it:} \emph{masīdu}, once more
  \item \emph{Say it:} say \emph{masīdu}
  \item \emph{Recall:} say the Telugu for a factory, then the Telugu for a jail, then say \emph{masīdu} again
\end{itemize}

\begin{tcolorbox}[breakable,colback=teal!4,colframe=teal!35!black,title={Before you move on}]
What does \te{మసీదు} mean? (\textquotedblleft{}A mosque\textquotedblright{}.) Say it once more.
\end{tcolorbox}
\section[\texorpdfstring{\te{చర్చి} (carci)}{చర్చి (carci)}]{\te{చర్చి} (carci) — a church}

\label{lesson:TE-C293-carci}



\begin{quote}
Before the new one: say the Telugu for a jail, then the Telugu for a mosque.
\end{quote}

\subsection*{You'll want to know: \te{చర్చి}}
\textbf{\te{చర్చి}} — \emph{carci} — \textquotedblleft{}a church\textquotedblright{}.

\begin{grammarlens}[title={What you've built}]
One more word for this chapter.
\end{grammarlens}

\subsection*{Your turn}
\begin{itemize}
\raggedright
  \item \emph{Say it:} \emph{carci}
  \item \emph{Say it:} \emph{carci}, once more
  \item \emph{Say it:} say \emph{carci}
  \item \emph{Recall:} say the Telugu for a jail, then the Telugu for a mosque, then say \emph{carci} again
\end{itemize}

\begin{tcolorbox}[breakable,colback=teal!4,colframe=teal!35!black,title={Before you move on}]
What does \te{చర్చి} mean? (\textquotedblleft{}A church\textquotedblright{}.) Say it once more.
\end{tcolorbox}
\section[\texorpdfstring{\te{పంచాయతీ} (pañcāyatī)}{పంచాయతీ (pañcāyatī)}]{\te{పంచాయతీ} (pañcāyatī) — the village council}

\label{lesson:TE-C293-pancayati}



\begin{quote}
Before the new one: say the Telugu for a mosque, then the Telugu for a church.
\end{quote}

\subsection*{You'll want to know: \te{పంచాయతీ}}
\textbf{\te{పంచాయతీ}} — \emph{pañcāyatī} — \textquotedblleft{}the village council\textquotedblright{}.

\begin{grammarlens}[title={What you've built}]
One more word for this chapter.
\end{grammarlens}

\subsection*{Your turn}
\begin{itemize}
\raggedright
  \item \emph{Say it:} \emph{pañcāyatī}
  \item \emph{Say it:} \emph{pañcāyatī}, once more
  \item \emph{Say it:} say \emph{pañcāyatī}
  \item \emph{Recall:} say the Telugu for a mosque, then the Telugu for a church, then say \emph{pañcāyatī} again
\end{itemize}

\begin{tcolorbox}[breakable,colback=teal!4,colframe=teal!35!black,title={Before you move on}]
What does \te{పంచాయతీ} mean? (\textquotedblleft{}The village council\textquotedblright{}.) Say it once more.
\end{tcolorbox}
\section[\texorpdfstring{\te{అటక} (aṭaka)}{అటక (aṭaka)}]{\te{అటక} (aṭaka) — a loft}

\label{lesson:TE-C293-ataka}



\begin{quote}
Before the new one: say the Telugu for a church, then the Telugu for the village council.
\end{quote}

\subsection*{You'll want to know: \te{అటక}}
\textbf{\te{అటక}} — \emph{aṭaka} — \textquotedblleft{}a loft\textquotedblright{}.

\begin{grammarlens}[title={What you've built}]
One more word for this chapter.
\end{grammarlens}

\subsection*{Your turn}
\begin{itemize}
\raggedright
  \item \emph{Say it:} \emph{aṭaka}
  \item \emph{Say it:} \emph{aṭaka}, once more
  \item \emph{Say it:} say \emph{aṭaka}
  \item \emph{Recall:} say the Telugu for a church, then the Telugu for the village council, then say \emph{aṭaka} again
\end{itemize}

\begin{tcolorbox}[breakable,colback=teal!4,colframe=teal!35!black,title={Before you move on}]
What does \te{అటక} mean? (\textquotedblleft{}A loft\textquotedblright{}.) Say it once more.
\end{tcolorbox}
